\section{Related Work}\label{sec:relwork}


A cornucopia of IRs has been presented in the literature to better expose desirable
program properties for optimizations. For brevity, we restrict our discussion to the
most prominent IRs, only highlighting their strengths and weaknesses in comparison to
the RVSDG, and refer the reader to Stanier~\etal~\cite{Stanier13} for a greater survey.

\subsection{Control (Data) Flow Graph}
The Control Flow Graph (CFG)~\cite{Allen70} exposes the intra-procedural control flow of a
function. Its nodes represent basic blocks, \ie, an ordered list of operations without
branches or branch targets, and its edges represent the possible control flow paths between these
nodes. This explicit exposure of control flow simplifies certain analyses, such as loop
identification or irreducibility detection, and enables simple target code generation. The CFG's
translation to SSA form~\cite{Cytron91}, or one of its variants, such as gated SSA~\cite{Tu95},
thinned gated SSA~\cite{Havlak93}, or future gated SSA~\cite{Ding14}, additionally improves the
efficiency of data flow optimizations~\cite{Wegman91,Rosen88}. These properties along with its
simple construction from a language's abstract syntax tree made the CFG in SSA form the predominant
IR for imperative language compilers~\cite{Stanier13}, such as LLVM~\cite{Lattner04} and
GCC~\cite{gcc}. However, the CFG has also been criticized as an IR for optimizing
compilers~\cite{Ferrante87,Johnson03,Johnson04,Lawrence07,Weise94,Zaidi15,Zaidi15-2}:

\begin{enumerate}
	\item It is incapable of representing inter-procedural information. It requires additional IRs,
		\eg, the call graph, to represent such information.

	\item It provides no structural information about a procedure's body. Important structures, such
		as loops, and their nesting needs to be constantly (re-)discovered for optimizations, as well
		as normalized to make them amenable for transformations.

	\item It emphasizes control dependencies, even though many optimizations are based on the flow
		of data. This is somewhat mitigated by translating it to SSA form or one of its variants, but
		in turn requires SSA restoration passes~\cite{Choi96} to ensure SSA invariants.

	\item It is an inherently sequential IR. The operations in basic blocks are listed in a
		sequential order, even if they are not dependent on each other. Moreover, this
		sequentialization also exists for structures such as loops, as two independent loops can only
		be represented in sequential order. Thus, the CFG is by design incapable of explicitly
		encoding independent operations.

	\item It provides no means to encode additional dependencies other than control and true data
		dependencies. Other information, such as loop-carried dependencies or alias information, must
		regularly be recomputed and/or memoized in addition to the CFG.
\end{enumerate}

The Control Data Flow Graph (CDFG)~\cite{Namballa04} tries to mitigate the sequential nature of
the CFG by replacing the sequence of operations in basic blocks with the Data Flow
Graph (DFG)~\cite{Dennis80}, an acyclic graph that represents the flow of data between operations.
This relaxes the strict ordering within a basic block, but does not expose instruction level
parallelism beyond basic block boundaries or between program structures.

\subsection{Program Dependence Graph/Web}
The Program Dependence Graph (PDG) \cite{Ferrante87,Horwitz288} combines control and data flow
within a single representation. It features data and control flow edges, as well as statement,
predicate, and region nodes. Statement nodes represent operations, predicate nodes represent
conditional choices, and region nodes group nodes with the same control dependency. If a
region's control dependencies are fulfilled, then its children can be executed in
parallel. Horwitz \etal~\cite{Horwitz88} extended the PDG to model inter-procedural dependencies
by incorporating procedures into the graph.

The PDG improves upon the CFG by employing region nodes to relax the overly restrictive sequence
of operations. This relaxed sequence combined with the unified representation of data and control
dependencies simplifies complex optimizations, such as code vectorization~\cite{Baxter89} or the
extraction of thread-level parallelism~\cite{Ottoni05,Sarkar91}. However, the unified data and
control flow representation results in a large number of edge types, five in Ferrante
\etal~\cite{Ferrante87} and four in Horwitz \etal~\cite{Horwitz288}, which need to be maintained
to ensure the graph's invariants. The PDG suffers from aliasing and side-effect problems, as it
supports no clear distinction between data held in register and memory. This complicates or can
even preclude its construction altogether~\cite{Johnson04}. Moreover, program structure and SSA
form still need to be discovered and maintained.

The Program Dependence Web (PDW)~\cite{Ottenstein90} extends the PDG and gated SSA~\cite{Tu95}
to provide a unified representation for the interpretation of programs using control-, data-, or
demand-driven execution models. This simplifies the mapping of programs written in different
paradigms, such as the imperative or functional paradigm, to different architectures, such as
Von-Neumann and dataflow architectures. In addition to the elements of the PDG, the PDW adds $\mu$
nodes to manage initial and loop-carried values and $\eta$ nodes to manage loop-exit
values. Campbell \etal~\cite{Campbell93} further refined the definition of the PDW by replacing
$\mu$ nodes with $\beta$ nodes and eliminating $\eta$ nodes. As the PDW is based on the PDG, it
suffers from the same aliasing and side-effect problems. PDW's additional constructs
further complicate graph maintenance and its construction is elaborate, requiring three
additional passes over a PDG, and is limited to programs with reducible control flow.



\subsection{Value (State) Dependence Graph}
The Value Dependence Graph (VDG)~\cite{Weise94} abandons the explicit representation of control
flow and only models the flow of values using ports. Its nodes represent simple operations, the
selection between values, or functions, using recursive functions to model loops.
The VDG is implicitly in SSA form and abandons the sequential order of operations from the CFG, as
each node is only dependent on its values. However, modeling only data flow between stateful
computations raises a significant problem in terms of preservation of program semantics, as the
"evaluation of the VDG may terminate even if the original program would not..."~\cite{Weise94}.


The Value State Dependence Graph (VSDG)~\cite{Johnson03,Johnson04} addresses the VDG's termination
problem by introducing state edges. These edges are used to model the sequential
execution of stateful computations. In addition to nodes for representing simple operations
and selection, it introduces nodes to explicitly represent loops. Like the VDG, the VSDG is
implicitly in SSA form, and nodes are solely dependent on required operands, avoiding a sequential
order of operations. However, the VSDG supports no inter-procedural constructs, and its selection
operator is only capable of selecting between two values based on a predicate. This complicates
destruction, as selection nodes must be combined to express conditionals. Even worse, the VSDG
represents all nodes as a flat graph, which simplifies optimizations~\cite{Johnson04}, but has a
severe effect on evaluation semantics. Operations with side-effects are no longer guarded by
predicates, and care must be taken to avoid duplicated evaluation of these operations. In fact,
for graphs with stateful computations, lazy evaluation is the only safe strategy~\cite{Lawrence07}.
The restoration of a program with an eager evaluation semantics complicates destruction immensely,
and requires a detour over the PDG to arrive at a unique CFG~\cite{Lawrence07}.
Zaidi \etal~\cite{Zaidi15-2,Zaidi15} adapted the VSDG to spatial hardware and sidestepped this
problem by introducing a predication-based eager/dataflow semantics. The idea is to effectively
enforce correct evaluation of operations with side-effects by using predication. While this seems
to circumvent the problem for spatial hardware, it is unclear what the performance implications
would be for conventional~processors.

The RVSDG solves the VSDG's eager evaluation problem by introducing regions. %
These regions enable the modeling of control flow constructs as nested nodes, and the guarding of
operations with side-effects. This avoids any possibility of duplicated evaluation, and in turn
simplifies RVSDG destruction. Moreover, nested nodes permit the explicit encoding of a program's
hierarchical structure into the graph, further simplifying optimizations.


